\subsection{The degree of the exceptional row sublattice}
\label{rank:exceptional-degree}

Fix a translate and use \(d,n,\rho,D_f,s,U,J\) of
Propositions~\ref{rank:transfer} and~\ref{rank:wide}. Among the target
HN cutoffs with positive codimension \(b<2nD_f\), take the one of
largest codimension \(b\). Call its saturated lattice \(W_*\).
If there is no such cutoff put \(W_*=R,\ b=0\). Give
\(A=(R/W_*)^\vee\) its dual quotient metric. It is the saturated
sublattice of the primal row lattice annihilating \(W_*\).
Write \(\widehat{\deg}A=b\mu(A)\) in the paper's squared-norm
normalization and set the degree of the zero lattice to zero.

\begin{proposition}[A signed degree at one cutoff]\label{rank:signed-degree}
With \(\tau=\rho/10\), the actual full-frame and column hypotheses of
Proposition~\ref{rank:wide} imply
\[
 s-\mu(R)\ \le\ \frac{J+(1-\tau)U}{\tau}
              +\frac{b(s+U)+\widehat{\deg}A}{d}.
\]
The last summand is signed. It can be averaged over the translates
before comparison with the arithmetic budget.
\end{proposition}
\begin{proof}
Write the centred target slopes as \(t_1\le\cdots\le t_d\) with
multiplicity. Replace \(t_1,\ldots,t_b\) by \(U\); denote the resulting
rank law by \(T'\). For \(v\le U\), if the number of slopes below
\(v\) is at most \(b\), then \(\Pr(T'<v)=0\). Otherwise the cutoff
has codimension at least \(2nD_f\), by the maximal choice of \(b\).
Proposition~\ref{rank:wide} and the integral filtration comparison give
\(\tau\Pr(T'<v)\le\tau\Pr(T<v)\le\Pr(S'<v)\) for the repaired
source law. For \(v>U\) the source probability is one. Integrating
as in Proposition~\ref{rank:transfer} gives
\(-\mathbb E T'\le (J+(1-\tau)U)/\tau\).
The HN quotient has slopes \(s+t_1,\ldots,s+t_b\).
Degree changes sign on taking its Hermitian dual by the determinant
pairing and the inverse dual norm. Hence
\[
 \mathbb ET'-\mathbb ET
      =\frac{bU-\sum_{j\le b}t_j}{d}
      =\frac{b(s+U)+\widehat{\deg}A}{d}.
\]
This includes \(b=d\) and coincident HN slopes, since an HN cutoff
only occurs at the end of a slope block. The finite rank-law argument
is also proved in \code{ExceptionalRowDegree.lean}.
\end{proof}

There is a direct arithmetic formula for the last degree. For \(b>0\)
take any rational map \(P\) from \(A_K\) to a Hermitian space with
lattice \(M\), injective on \(A_K\). Put \(L=M\cap P(A_K)\) with
induced metrics. At a finite place \(w\), let \(m_w\) be the
valuation of \(\det P:\det A_w\to\det L_w\) in generators of the
two local determinant lattices. It is the sum of the signed Smith
exponents, and vanishes at all but finitely many places. At \(v\mid
\infty\) write \(j_v=\|\det(P|_{A_v})\|\) using precisely these
determinant metrics. Since \(K\) is totally complex, write
\(\langle f\rangle_\infty=(2/[K:\mathbf Q])\sum_{v\mid\infty}f_v\).
Then
\[
 \widehat{\deg}A=\widehat{\deg}L+
      \big\langle \log j_v^2\big\rangle_\infty
          -2\sum_{w\nmid\infty}
                 \frac{\log Nw}{[K:\mathbf Q]}\,m_w .
\]
Indeed, evaluate both degrees on a nonzero rational generator of
\(\det A_K\) and its image. At a finite place its order increases
by \(m_w\), and at infinity its squared norm is multiplied by
\(j_v^2\). Changing the generator cancels by the product formula.
In particular one upper bound uses the height of the HN polygon of
\(M\) at rank \(b\) for \(\widehat{\deg}L\). This follows from the
definition of that polygon as the concave majorant of sublattice
degrees. See \href{https://arxiv.org/abs/0803.2583}{H.~Chen,
\emph{Positive degree and arithmetic bigness}, \S2}.
The Jacobian uses just the \(b\) directions in \(A\). Neither
necessity nor a cost for the remaining row directions enters this
formula. If a map is initially given on another level of the tower,
the equality requires its transport to the actual \(A\) with the
specified finite lattices and metrics.

For comparison with the original scale, put \(X=x^{C_H}\ge20\).
The identity-source specialization in Proposition~\ref{rank:wide}
gives \(\eta\le5/X\) for the full maps. The original size bounds
give \(\alpha\le d_e\le1/(9X^2)\) and
\(e\le e_0\le e_B/2\le1/(32X^2)\). Thus
\[
    D_f\le \frac5X+\frac{49}{72X^2},\qquad
    \frac bd\le \frac{11}{\rho X}.
\]
If the signed degree of \(A\) satisfies \(bs+\widehat{\deg}A\le
bXb_0\), where \(b_0\) denotes the original scale parameter
\(\mathtt{sc.b}\), its whole contribution is at most
\(11b_0/\rho + bU/d\). This bound asks only for the average slope
on the exceptional quotient, even if an individual slope lies
arbitrarily far below it. The original floor provides it. An
independent arithmetic upper bound on the displayed determinant
for the actual row sublattice would serve the same purpose.

\paragraph{The cost of a missed bottom slope.}
Let \(\lambda_1\le\cdots\le\lambda_d\) be the target slopes with
multiplicity. Write \(b_*\) for the largest positive codimension of
an HN cutoff below \(2nD_f\), taking \(b_*=0\) if there is none.
Suppose \(0<b_*<d\), and put \(c=\lambda_{b_*+1}\).
Under the frame and column assumptions above, with \(\tau=\rho/10\),
\[
 s-\mu(R)\ \le\ \frac{J+(1-\tau)U}{\tau}
             +\frac1d\sum_{j=1}^{b_*}(c-\lambda_j).
\]
Indeed replace each \(\lambda_j\) by \(\max(c,\lambda_j)\). Below
and at \(c\) the strict lower tail of this list is empty. Above
\(c\) its lower tail agrees with the original one, and any proper
cutoff there has codimension greater than \(b_*\). The whole-space
and zero-space cutoffs satisfy the same comparison whenever they
occur, using the large-quotient bound for codimension \(d\).
Apply the distribution inequality in Proposition~\ref{rank:transfer}
to the raised list and subtract the change of mean. For \(b_*=1\)
the exact additional allowance is \((\lambda_2-\lambda_1)/d\).

\paragraph{Replication of the extreme slope on the marked row space.}
Use the principal marked row realization from the original passage
between levels. Write \(M=R^\vee\) for the primal sharp row lattice
and \(M_G\) for the common lattice, with finite group \(G\) acting
integrally and unitarily at every infinite place. Let \(P\) be the
sharp character subgroup, \(h=|P|\), and let \(e\) be its character
projector. The sharp comparison supplies a rational identification
\(\jmath:M_K\simeq e(M_G)_K\). For \(L=M_G\cap e(M_G)_K\) it gives
\[
             hL\subseteq\jmath M\subseteq L .
\]
The two metrics differ by a common scalar at each infinite place,
so conjugation of an endomorphism by \(\jmath\) preserves its
operator norm. The original sharp row tie identifies this metric
and integral lattice with the ones defining the target dual.
\src{src/Formalism/LocalToGlobal/GlobalTheory/DivPassage.lean}{70}{metric and lattice conditions},
\src{src/Formalism/LocalToGlobal/IndefiniteReal/CutPassage3SOfSteps.lean}{71}{unitarity, stability and the lattice inclusions} and
\src{src/Formalism/LocalToGlobal/Divisibility/RowSideTie.lean}{259}{identification of the sharp rows}.

Decompose \((M_G)_K\) into copies of the occurring irreducible
\(K[G]\)-modules \(V_s\). Suppose every occurring \(eV_s\ne0\)
has dimension at least \(n_{\min}\). For each \(g\in G\) the map
\[
               T_g=h^2\jmath^{-1}e g e\jmath
\]
preserves \(M\). Indeed \(g\) preserves \(M_G\), \(he\) has
integral group-ring coefficients, and the remaining factor \(h\)
puts the image of \(L\) back in \(\jmath M\). The orthogonal
projector \(e\) and unitarity give \(\|T_g\|_v\le h^2\) at every
infinite place. Its squared-norm slope cost is at most \(4\log h\).

Let \(\mu_1\) be the largest slope of \(M\). The whole highest HN
step has minimum slope \(\mu_1\). Its images under all \(T_g\) lie
in the HN step at \(\mu_1-4\log h\), by the integral image
inequality. Pick a nonzero vector in the generic fibre of the
highest step. Its projection to some irreducible copy \(V_s\)
is a nonzero \(v\in eV_s\). Irreducibility gives
\(\operatorname{span}_K\{gv:g\in G\}=V_s\). Thus the projected
images of this vector under the \(T_g\) span \(eV_s\), and the
lower step has rank at least \(n_{\min}\). Duality yields
\[
             \lambda_{n_{\min}}-\lambda_1\le4\log h .
\]
This argument uses the images of one HN step under the compressed
finite-group operators. It requires no invariance of the word periods
under those operators.
\src{src/Formalism/HarderNarasimhan/FinalCor.lean}{9}{image in an HN cutoff}
and \src{src/Formalism/HarderNarasimhan/HarderNarasimhanProofs/SlopeIneqMean.lean}{128}{the squared operator norm}.
The usual height form of the filtered image inequality is also
given by \href{https://arxiv.org/abs/0803.2583}{H.~Chen,
\emph{Positive degree and arithmetic bigness}, \S2.5}.

In the original marked realization the primitive row and finite
character hypotheses give
\(n_s=\dim eV_s\ge(1-d_e)(g_0g_1)^{fk}\), with
\(d_e=8fr/\min(g_0,g_1)\). In particular \(n_{\min}>1\) on the
late stages under the original size bounds. Since
\(\log h\le\log|G|=\lambda_0\le4b_0\), the extra mean allowance
when \(b_*<n_{\min}\) is at most \(16(b_*/d)b_0\).
It is \(o(E_x)\) on the selected sizes, as \(E_x\ge n_Tmb_0 x^\beta\).
In particular an isolated one-dimensional exceptional cutoff costs
at most \(16b_0/d\), including when every word misses the line.
The finite-group realization and its metric and integral ties are
hypotheses of this deduction. A bound on the number of row labels
alone does not supply them. For the entire window the possible
cutoffs with \(b_*\ge n_{\min}\) still require an estimate.
\src{src/Development/AbelianVarietyTheory/Proofs/DivisorProofs/PassIn3OfRows.lean}{23}{marked character multiplicities}
and \src{src/Development/AbelianVarietyTheory/Proofs/DivisorProofs/IndefiniteScaleThetaDegree.lean}{12}{hypotheses at the selected realization};
\src{src/Formalism/SharpRestriction/Divisibility/HookProducts/LambdaZeroLeFourB.lean}{11}{the logarithmic group-size bound}.

The same operators reduce the remaining tests to subspaces stable
under \(eK[G]e\). At a centred target threshold \(v\), let \(A_v\)
be the dual of the quotient formed by slopes strictly below \(s+v\).
In primal coordinates it is the strict upper HN step at \(-s-v\).
Set \(B_v=\sum_{g\in G}T_g A_v\) and \(C=4\log h\). Then
\[
               A_v\subseteq B_v\subseteq A_{v+C}.
\]
The identity group element gives the first inclusion. The image
inequality applied to the finitely many slopes of the strict step
gives the second. The span of \(ege\) over \(g\in G\) is the
algebra \(eK[G]e\), so \(B_v\) is a submodule for this algebra.
Suppose the weighted rank comparison of Proposition~\ref{rank:transfer}
holds with ratio \(\tau\) on the quotients dual to the nonzero
\(B_v\). By inclusion and the source filtration inequality this gives
\(\tau\Pr(T<v)\le\Pr(S'<v+C)\) for the original centred slopes.
Applying the same mean comparison to \(T+C\) costs at most \(C\)
in the target mean, once for the entire flag. Every \(B_v\) of
dimension at least \(2nD_f\) already passes with the original rank
weights and ratio \(\rho/10\). The outstanding rank tests in this
version concern only the smaller nonzero submodules \(B_v\).

\paragraph{How many base-level slopes a larger cutoff requires.}
Return to both marked levels in the original row realization. Write
\(M_0,M_1\) for the primal base and sharp lattices, \(e_i,\jmath_i\)
for their projectors and embeddings, \(h_i=|P_i|\),
\(d_0=\operatorname{rk}M_0\), \(d=\operatorname{rk}M_1\).
In addition to the sandwich above, the original finite-place
averaging formula gives \(h_i e_i M_G\subseteq\jmath_i M_i\).
Indeed a principal lattice element has a local representative after
multiplication by some element outside the given prime ideal.
Averaging that representative by \(h_i e_i\) produces an integral
row at level \(i\), with the same denominator. Intersect these
local conditions to obtain the asserted inclusion. The metric
identity gives \(\|x\|_i=h_i^{-1/2}\|\jmath_i x\|_G\).
\src{src/Formalism/LocalToGlobal/IndefiniteReal/RowGSurj.lean}{159}{finite-place integral averaging}
and \src{src/Formalism/LocalToGlobal/IndefiniteReal/CxPassMet.lean}{160}{metric scaling}.
Thus for \(g\in G\) the integral map
\(h_i\jmath_i^{-1}e_i g\jmath_j:M_j\to M_i\) has norm at most
\(\sqrt{h_i h_j}\) at every infinite place.

Write \(\alpha_1\ge\cdots\ge\alpha_d\) and
\(\beta_1\ge\cdots\ge\beta_{d_0}\) for the sharp and base
primal slopes. Put \(m_1=\min_s n_s\), \(M_1'=\max_s n_s\)
over the occurring good types, and \(L_0=\log(h_0h_1)\).
The base projector has rank one on each such irreducible type,
and the sharp projector has rank \(n_s\). In the common space take
the \(K[G]\)-span of an HN step at threshold \(\theta\) at either
level. If this subrepresentation has multiplicities \(a_s\), its
base and sharp projections have dimensions \(\sum_s a_s\) and
\(\sum_s n_s a_s\). Its projection back to the starting level
contains the starting step. Its projection to the other level lies
in the HN step at \(\theta-L_0\) by the integral maps just
constructed. Applying this in both directions gives
\[
 \beta_{\lceil i/m_1\rceil}-L_0
      \ \le\ \alpha_i\ \le\
       \beta_{\lceil i/M_1'\rceil}+L_0 .
\]
The left bound holds for \(i\le m_1 d_0\), and the right bound
holds for all \(1\le i\le d\). At the sharp level alone the same
averaging improves the earlier slope cost to \(2\log h_1\).
In particular \(b_*<m_1\) costs at most
\(2(b_*/d)\log h_1\le 8(b_*/d)b_0\).

For \(b=b_*>0\) and \(j=1+\lfloor b/m_1\rfloor\le d_0\),
duality gives the following bound on the exact clipping allowance:
\[
 \frac{1}{d}\sum_{i=1}^{b}(\alpha_i-\alpha_{b+1})
 \le \frac{2L_0 b}{d}+
         \frac{M_1'}{d}\sum_{\ell<j}(\beta_\ell-\beta_j).
\]
Indeed \(\alpha_{b+1}\ge\beta_j-L_0\), and each positive difference
on the right can occur at most \(M_1'\) times in the sum of the
upper bounds. The multiplicity estimates give
\(M_1'/d\le(1+d_e)/((1-d_e)d_0)\).
Since \(b<2nD_f\), \(d=\rho n\), and
\(m_1d_0/d\ge(1-d_e)/(1+d_e)\), the required \(j\le d_0\)
holds on the late scales with \(\rho\ge1/2\), \(D_f=O(1/X)\)
and \(d_e=o(1)\). A larger missed cutoff therefore asks for the
top \(\lfloor b/m_1\rfloor\) base slopes relative to the next one.
The common marked representation supplies the multiplicities in
this inequality. The sizes of those base slope differences are
additional arithmetic data.

There is a single sufficient base estimate. For \(q=j-1\) and
\(\bar\beta=\mu(M_0)\), monotonicity gives
\(d_0\bar\beta\le q\beta_1+(d_0-q)\beta_{q+1}\). Hence
\[
 \frac{1}{d_0}\sum_{\ell=1}^{q}(\beta_\ell-\beta_{q+1})
     \le\frac{q}{d_0-q}\big(\beta_1-\bar\beta\big).
\]
Here \(q/d_0\le (b/d)(M_1'/m_1)=O(1/X)\). Consequently a bound
\(\mu_{\max}(M_0)-\mu(M_0)=O(X E_x)\) on the required translates
would pay for every residual cutoff at the scale \(E_x\), under
the other frame and column hypotheses. The original base-to-sharp
passage contains \(\mu(M_0^\vee)+H_0\) with an input
\(H_0\ge\mu_{\max}(M_0)\). A dual-mean upper bound alone controls
only the second term of \(\mu_{\max}(M_0)-\mu(M_0)\).
\src{src/Formalism/LocalToGlobal/IndefiniteForms/IndefSide.lean}{339}{the input to the original passage}.

Even a saturated quotient of a standard Euclidean lattice can
have an arbitrarily large gap. For instance let
\(a+bi=(2+i)^N\), so \(a^2+b^2=5^N\) and \(\gcd(a,b)=1\)
by unique factorization in \(\mathbf Z[i]\). In
\(\mathbf Z^3/\mathbf Z(a,b,0)\), with the quotient metric, a
Bézout complement in the first two coordinates has squared norm
\(5^{-N}\), and the third coordinate has squared norm one,
orthogonally. The gap is \(N\log 5\) in the squared-norm
convention. This example tests what follows from saturation and
the quotient metric alone. For the actual row lattice the answer
depends on the subspace removed by the period construction.

One can pin down the remaining metric as well. Pull the actual
base lattice at every finite place back through \(\iota_{\mu_0}\),
and on this same lattice use the rational quotient metric at
infinity. Let \(\beta^Q_\ell\) be its primal slopes. On each
occupied orthogonal type \(p\) the original squared row norm,
in these coordinates, is the quotient squared norm times
\(\kappa_v\prod_{j=1}^{m}\lambda_{j,p,v}\).
The base-type count expresses the rank of the type as a sum of
nonnegative row multiplicities. A positive summand supplies a
label with a nonzero isotypic component at every lift. The
original component identification therefore applies at each
embedding to every occupied type. With the original recorded
two-sided estimate
\[
 |\log\lambda_{j,p,v}-H_j|\le\Delta,\qquad
 \Delta=C f n_V^2\log(b_0\,\mathrm{rd}_E\,g_0g_1 n_V),
\]
the common scalar \(\kappa_v\exp(\sum_jH_j)\) shifts all slopes
by the same amount. After removing that scalar the two squared
metrics differ by factors between \(e^{-m\Delta}\) and
\(e^{m\Delta}\). The integral identity map in both directions
bounds the error in each slope by \(m\Delta\). Consequently
\[
 \left|\frac{1}{d_0}\sum_{\ell=1}^{q}
  \big[(\beta_\ell-\beta_{q+1})
         -(\beta^Q_\ell-\beta^Q_{q+1})\big]\right|
    \le \frac{2q}{d_0}m\Delta .
\]
On the original sizes \(m\Delta=x^{O(1)}=o(b_0)\).
Thus the common heat factors and these typewise errors fit in
the mean budget. What remains here is the slope spread of the
rational quotient with its actual glued finite lattices.
\src{src/Formalism/LocalToGlobal/IndefiniteForms/IndefSide.lean}{304}{the scalar metric on each type},
\src{src/Development/AbelianVarietyTheory/Proofs/IwasawaDecomposition/FinrankLiftLabelEqFinrankFace3.lean}{415}{the base type count},
\src{src/Development/AutomorphicArithmetic/ThetaCorrespondence/ThetaPairLifts/ThirdPlaceSeries/RecInPartsExists.lean}{30}{identification of the type factors in the recorded realization} and
\src{src/Formalism/SharpRestriction/Divisibility/LocalQuadraticParts.lean}{257}{the two-sided heat estimate}.
